
\documentclass[11pt]{article}

\usepackage{amsmath,amssymb,amsthm}
\usepackage{mathpartir}
\usepackage{geometry}
\usepackage{hyperref}

\geometry{margin=1in}

\title{Language Definition}
\author{}
\date{}

\begin{document}
\maketitle

%%%%%%%%%%%%%%%%%%%%%%%%%%%%%%%%%%%%%%%%%%%%%%%%%%%%%%%%%%%%
\section{Abstract Syntax}
%%%%%%%%%%%%%%%%%%%%%%%%%%%%%%%%%%%%%%%%%%%%%%%%%%%%%%%%%%%%

\subsection{Values}
\[
\begin{array}{rcl}
v \;\;::=& n \\
        &|& b \\
        &|& C(v_1,\dots,v_n) \\
\end{array}
\]

where $n \in \mathbb{Z}$, $b \in \{\texttt{true}, \texttt{false}\}$, and $C$ is a constructor name.

\subsection{Types}

\[
\begin{array}{rcl}
T \;\;::=& \texttt{int} \\
        &|& \texttt{bool} \\
        &|& V \\
\end{array}
\]


Variant types are labeled sums:

\[
V ::= \langle C_1 : \overline{T_1}, \dots, C_n : \overline{T_n} \rangle
\]

where each $\overline{T_i}$ denotes a (possibly empty) sequence of types and $C_i$ are
constructor names.

%%%%%%%%%%%%%%%%%%%%%%%%%%%%%%%%%%%%%%%%%%%%%%%%%%%%%%%%%%%%
\subsection{Expressions}
%%%%%%%%%%%%%%%%%%%%%%%%%%%%%%%%%%%%%%%%%%%%%%%%%%%%%%%%%%%%

\[
\begin{array}{rcl}
e \;\;::=& v
\\ &|& e_1 \;\texttt{op}\; e_2
\\ &|& \texttt{!}\; e
\\ &|& x
\end{array}
\]

where $\texttt{op} \in \{ +, -, *, /, \land, \lor, ==, <, > \}$ and $x$ is a variable name.

%%%%%%%%%%%%%%%%%%%%%%%%%%%%%%%%%%%%%%%%%%%%%%%%%%%%%%%%%%%%
\subsection{Statements}
%%%%%%%%%%%%%%%%%%%%%%%%%%%%%%%%%%%%%%%%%%%%%%%%%%%%%%%%%%%%

\[
\begin{array}{rcl}
s \;\;::=& \texttt{var } x : T
\\ &|& x := e
\\ &|& \texttt{if } e \;\{ s^* \}
\\ &|& \texttt{if } e \;\{ s^* \}\; \texttt{else } \{ s^* \}
\\ &|& \texttt{print } e
\\ &|& \texttt{variant } V = C_1(\overline{T_1}) \mid \dots \mid C_n(\overline{T_n})
\end{array}
\]

%%%%%%%%%%%%%%%%%%%%%%%%%%%%%%%%%%%%%%%%%%%%%%%%%%%%%%%%%%%%
\section{Typing Rules}
%%%%%%%%%%%%%%%%%%%%%%%%%%%%%%%%%%%%%%%%%%%%%%%%%%%%%%%%%%%%
\[ 
\Gamma \vdash e : T
\]
\[
\Delta \vdash C : (\overline{T}) : V
\]
\subsection{Base Types}

\begin{mathpar}
\inferrule*[right=T-Int]
{ }
{\Gamma \vdash n : \texttt{int}}

\inferrule*[right=T-Bool-True]
{ }
{\Gamma \vdash \texttt{true} : \texttt{bool}}

\inferrule*[right=T-Bool-False]
{ }
{\Gamma \vdash \texttt{false} : \texttt{bool}}
\end{mathpar}

%%%%%%%%%%%%%%%%%%%%%%%%%%%%%%%%%%%%%%%%%%%%%%%%%%%%%%%%%%%%
\subsection{Variables}

\begin{mathpar}
\inferrule*[right=T-Var-Use]
{x : T \in \Gamma}
{\Gamma \vdash x : T}
\end{mathpar}

%%%%%%%%%%%%%%%%%%%%%%%%%%%%%%%%%%%%%%%%%%%%%%%%%%%%%%%%%%%%
\subsection{Binary Operators}

\begin{mathpar}
\inferrule*[right=T-Arithmetic]
{\Gamma \vdash e_1 : \texttt{int} \\ \Gamma \vdash e_2 : \texttt{int}}
{\Gamma \vdash e_1 + e_2 : \texttt{int}}

\inferrule*[right=T-Logical]
{\Gamma \vdash e_1 : \texttt{bool} \\ \Gamma \vdash e_2 : \texttt{bool}}
{\Gamma \vdash e_1 \land e_2 : \texttt{bool}}

\inferrule*[right=T-Comparison]
{\Gamma \vdash e_1 : \texttt{int} \\ \Gamma \vdash e_2 : \texttt{int}}
{\Gamma \vdash e_1 \textless e_2 : \texttt{bool}}

\inferrule*[right=T-Equality]
{\Gamma \vdash e_1 : T \\ \Gamma \vdash e_2 : T}
{\Gamma \vdash e_1 == e_2 : \texttt{bool}}
\end{mathpar}

%%%%%%%%%%%%%%%%%%%%%%%%%%%%%%%%%%%%%%%%%%%%%%%%%%%%%%%%%%%%
\subsection{Unary Operators}

\begin{mathpar}
\inferrule*[right=T-Not]
{\Gamma \vdash e : \texttt{bool}}
{\Gamma \vdash \texttt{not } e : \texttt{bool}}
\end{mathpar}

%%%%%%%%%%%%%%%%%%%%%%%%%%%%%%%%%%%%%%%%%%%%%%%%%%%%%%%%%%%%
\subsection{Variants}

\begin{mathpar}
\inferrule*[right=T-Variant-Constructor]
{
\Gamma \vdash e_1 : T_1 ,\cdots, e_n : T_n \\ 
\Delta \vdash C : (T_1,\dots,T_n) : V
}
{
\Gamma, \Delta \vdash C(e_1,\dots,e_n) : V
}
\end{mathpar}

The constructor $C$ must be declared as part of variant $V$ with argument
types $(T_1,\dots,T_n)$.

%%%%%%%%%%%%%%%%%%%%%%%%%%%%%%%%%%%%%%%%%%%%%%%%%%%%%%%%%%%%
\section{Operational Semantics}
%%%%%%%%%%%%%%%%%%%%%%%%%%%%%%%%%%%%%%%%%%%%%%%%%%%%%%%%%%%%

The operational semantics are given as a big-step evaluation relation:
\[
\langle e, \sigma \rangle \Downarrow v
\]

%%%%%%%%%%%%%%%%%%%%%%%%%%%%%%%%%%%%%%%%%%%%%%%%%%%%%%%%%%%%
\subsection{Base Values}

\begin{mathpar}
\inferrule*[right=E-Int]
{ }
{\langle n, \sigma \rangle \Downarrow n}

\inferrule*[right=E-Bool]
{ }
{\langle \texttt{b}, \sigma \rangle \Downarrow \texttt{b}}
\end{mathpar}

%%%%%%%%%%%%%%%%%%%%%%%%%%%%%%%%%%%%%%%%%%%%%%%%%%%%%%%%%%%%
\subsection{Variables}

\begin{mathpar}
\inferrule*[right=E-Var]
{\sigma(x) = v}
{\langle x, \sigma \rangle \Downarrow v}
\end{mathpar}

%%%%%%%%%%%%%%%%%%%%%%%%%%%%%%%%%%%%%%%%%%%%%%%%%%%%%%%%%%%%
\subsection{Operations}

\begin{mathpar}
\inferrule
{
\langle e_1, \sigma \rangle \Downarrow n_1 \\
\langle e_2, \sigma \rangle \Downarrow n_2
}
{
\langle e_1 + e_2, \sigma \rangle \Downarrow n_1 + n_2
}
\end{mathpar}

\begin{mathpar}
\inferrule
{
\langle e_1, \sigma \rangle \Downarrow b_1 \\
\langle e_2, \sigma \rangle \Downarrow b_2
}
{
\langle e_1 \land e_2, \sigma \rangle \Downarrow b_1 \land b_2
}
\end{mathpar}

\begin{mathpar}
\inferrule
{
\langle e, \sigma \rangle \Downarrow b
}
{
\langle \texttt{!} e, \sigma \rangle \Downarrow \neg b
}
\end{mathpar}

\begin{mathpar}
\inferrule
{
\langle e_1, \sigma \rangle \Downarrow n_1 \\
\langle e_2, \sigma \rangle \Downarrow n_2
}
{
\langle e_1 \textless e_2, \sigma \rangle \Downarrow (n_1 < n_2)
}
\end{mathpar}

\begin{mathpar}
\inferrule
{
\langle e_1, \sigma \rangle \Downarrow v_1 \\
\langle e_2, \sigma \rangle \Downarrow v_2
}
{
\langle e_1 == e_2, \sigma \rangle \Downarrow (v_1 == v_2)
}
\end{mathpar}
%%%%%%%%%%%%%%%%%%%%%%%%%%%%%%%%%%%%%%%%%%%%%%%%%%%%%%%%%%%%
\subsection{Variants}

\begin{mathpar}
\inferrule
{
\langle e_1, \sigma \rangle \Downarrow v_1 \\
\cdots \\
\langle e_n, \sigma \rangle \Downarrow v_n
}
{
\langle C(e_1,\dots,e_n), \sigma \rangle \Downarrow C(v_1,\dots,v_n)
}
\end{mathpar}

%%%%%%%%%%%%%%%%%%%%%%%%%%%%%%%%%%%%%%%%%%%%%%%%%%%%%%%%%%%%
\section{Statements}

Statement execution is defined as:
\[
\langle s, \sigma \rangle \Downarrow \sigma'
\]

\begin{mathpar}
\inferrule
{\langle e, \sigma \rangle \Downarrow v}
{\langle x := e, \sigma \rangle \Downarrow \sigma[x \mapsto v]}
\end{mathpar}

\begin{mathpar}
\inferrule
{\langle e, \sigma \rangle \Downarrow \texttt{true}}
{\langle \texttt{if } e \{ s^* \}, \sigma \rangle \Downarrow \sigma'}
\end{mathpar}

%%%%%%%%%%%%%%%%%%%%%%%%%%%%%%%%%%%%%%%%%%%%%%%%%%%%%%%%%%%%
\section{Conclusion}

This document formally specifies the abstract syntax, typing rules, and
operational semantics of the language as implemented in the provided C++
source files.

\end{document}
